% 分野別類題: 行列の対角化・固有値・固有ベクトル 解答例
% コンパイル: TEXINPUTS="../../../00_common:$TEXINPUTS" latexmk -xelatex answers.tex
\documentclass[a4paper,11pt]{article}
\usepackage{preamble}
\usepackage{shoakubox}

\begin{document}

\kamokumei{類題演習 解答例 --- 行列の対角化・固有値・固有ベクトル}

\daimon{【解法】固有値・固有ベクトルと対角化の解き方と導出}

\noindent
$n$ 次正方行列 $A$ に対し、
\[
A\bm{x} = \lambda \bm{x}, \qquad \bm{x} \neq \bm{0}
\]
を満たすスカラー $\lambda$ を $A$ の\textbf{固有値}、$\bm{x}$ を対応する\textbf{固有ベクトル}という。この式は $(A - \lambda I)\bm{x} = \bm{0}$ と同値であり、$\bm{x} \neq \bm{0}$ が存在するためには係数行列 $A - \lambda I$ が正則でないこと、すなわち
\[
\det(A - \lambda I) = 0
\]
（固有方程式・特性方程式）が必要十分である。これを $\lambda$ について解いて固有値を求め、各固有値 $\lambda_i$ について $(A - \lambda_i I)\bm{x} = \bm{0}$ の解空間（固有空間）から固有ベクトルを求める。

\medskip
\noindent
\textbf{対角化.} $A$ が $n$ 個の一次独立な固有ベクトル $\bm{p}_1, \dots, \bm{p}_n$（固有値 $\lambda_1, \dots, \lambda_n$）を持つとき、これらを列に並べた行列 $P = (\bm{p}_1 \ \cdots \ \bm{p}_n)$ は正則で、
\[
AP = (A\bm{p}_1 \ \cdots \ A\bm{p}_n) = (\lambda_1 \bm{p}_1 \ \cdots \ \lambda_n \bm{p}_n) = P D, \qquad D = \operatorname{diag}(\lambda_1, \dots, \lambda_n)
\]
が成り立つ。両辺に左から $P^{-1}$ を掛けて
\[
P^{-1} A P = D
\]
を得る。これが対角化である。相異なる固有値に対応する固有ベクトルは自動的に一次独立になるので、固有値がすべて相異なれば対角化可能である。

\medskip
\noindent
\textbf{対称行列の直交対角化（スペクトル定理）.} $A$ が実対称行列（$A^{\mathrm{T}} = A$）のとき、固有値はすべて実数であり、相異なる固有値に対応する固有ベクトルは互いに直交する。実際、$A\bm{x} = \lambda \bm{x}$, $A\bm{y} = \mu \bm{y}$（$\lambda \neq \mu$）ならば
\[
\lambda\, \bm{x}^{\mathrm{T}}\bm{y} = (A\bm{x})^{\mathrm{T}}\bm{y} = \bm{x}^{\mathrm{T}} A^{\mathrm{T}} \bm{y} = \bm{x}^{\mathrm{T}} A \bm{y} = \mu\, \bm{x}^{\mathrm{T}}\bm{y}
\]
より $(\lambda - \mu)\bm{x}^{\mathrm{T}}\bm{y} = 0$ であり $\lambda \neq \mu$ だから $\bm{x}^{\mathrm{T}}\bm{y} = 0$。同じ固有値に重複度がある場合もグラム・シュミットの直交化により固有空間内で正規直交基底を選べる。したがって各固有ベクトルを正規化して並べた行列 $P$ は直交行列（$P^{\mathrm{T}}P = I$、すなわち $P^{-1} = P^{\mathrm{T}}$）となり、
\[
P^{\mathrm{T}} A P = D
\]
という直交対角化が得られる。

\medskip
\noindent
\textbf{対角化を用いた $A^{n}$ の計算.} $P^{-1}AP = D$ のとき $A = PDP^{-1}$ であるから
\[
A^{n} = (PDP^{-1})^{n} = P D^{n} P^{-1}
\]
（途中の $P^{-1}P = I$ がすべて相殺するため）。$D$ は対角行列なので $D^{n} = \operatorname{diag}(\lambda_1^{n}, \dots, \lambda_n^{n})$ と容易に計算でき、$A^{n}$ を閉じた形で求められる。

\clearpage
\begin{mondaiboxS}{問1}
$A = \begin{pmatrix} 3 & 1 \\ 1 & 3 \end{pmatrix}$ の固有値・固有ベクトルを求め、直交行列 $P$ を用いて直交対角化 $P^{\mathrm{T}}AP = D$ を行え。
\end{mondaiboxS}
\kaitou
固有方程式は
\[
\det(A - \lambda I) = (3-\lambda)^{2} - 1 = \lambda^{2} - 6\lambda + 8 = 0
\]
より $\lambda = 4, 2$。

$\lambda = 4$: $A - 4I = \begin{pmatrix} -1 & 1 \\ 1 & -1 \end{pmatrix}$ より $x_1 = x_2$、固有ベクトルは $\bm{p}_1 = \dfrac{1}{\sqrt{2}}\begin{pmatrix}1\\1\end{pmatrix}$（正規化済み）。

$\lambda = 2$: $A - 2I = \begin{pmatrix} 1 & 1 \\ 1 & 1 \end{pmatrix}$ より $x_1 = -x_2$、固有ベクトルは $\bm{p}_2 = \dfrac{1}{\sqrt{2}}\begin{pmatrix}1\\-1\end{pmatrix}$。

$\bm{p}_1 \perp \bm{p}_2$（対称行列の異なる固有値の固有ベクトルは直交する）なので、そのまま直交行列
\[
P = \frac{1}{\sqrt{2}}\begin{pmatrix} 1 & 1 \\ 1 & -1 \end{pmatrix}
\]
が取れ、
\[
\boxed{\; P^{\mathrm{T}} A P = \begin{pmatrix} 4 & 0 \\ 0 & 2 \end{pmatrix} \;}
\]
（検算: $AP = \dfrac{1}{\sqrt2}\begin{pmatrix}3+1 & 3-1\\1+3 & 1-3\end{pmatrix} = \dfrac{1}{\sqrt2}\begin{pmatrix}4 & 2\\4 & -2\end{pmatrix}$、一方 $PD = \dfrac{1}{\sqrt2}\begin{pmatrix}1\cdot4 & 1\cdot2\\1\cdot4 & -1\cdot2\end{pmatrix} = \dfrac{1}{\sqrt2}\begin{pmatrix}4&2\\4&-2\end{pmatrix}$ で一致し、また $P^{\mathrm{T}}P = \dfrac12\begin{pmatrix}1+1 & 1-1\\1-1&1+1\end{pmatrix} = I$ も確認できる。）

\clearpage
\begin{mondaiboxS}{問2}
$A = \begin{pmatrix} 2 & 1 & 1 \\ 1 & 2 & 1 \\ 1 & 1 & 2 \end{pmatrix}$ の固有値・固有ベクトルを求め、直交行列 $P$ を用いて直交対角化 $P^{\mathrm{T}}AP = D$ を行え。
\end{mondaiboxS}
\kaitou
$A = I + J$（$J$ はすべての成分が $1$ の $3\times3$ 行列）と書ける。$J$ の固有値は、階数 $1$・trace $3$ より $3, 0, 0$（$\lambda=3$ の固有ベクトルは $(1,1,1)^{\mathrm{T}}$、$\lambda=0$ の固有空間は $x+y+z=0$ を満たす平面）であるから、$A = I+J$ の固有値は
\[
\lambda = 1+3 = 4 \ (\text{単根}), \qquad \lambda = 1+0 = 1 \ (\text{重根})
\]
実際に検算すると、$\det(A - \lambda I)$ を直接展開しても
\[
\det(A-\lambda I) = -(\lambda-4)(\lambda-1)^{2}
\]
となり一致する。

$\lambda = 4$: 固有空間は $(1,1,1)^{\mathrm{T}}$ 方向。正規化して $\bm{p}_1 = \dfrac{1}{\sqrt{3}}(1,1,1)^{\mathrm{T}}$。

$\lambda = 1$: 固有空間は平面 $x+y+z=0$。この中で正規直交基底を（グラム・シュミットで）選ぶと、例えば
\[
\bm{p}_2 = \frac{1}{\sqrt{2}}(1,-1,0)^{\mathrm{T}}, \qquad
\bm{p}_3 = \frac{1}{\sqrt{6}}(1,1,-2)^{\mathrm{T}}
\]
（$\bm{p}_2 \cdot \bm{p}_3 = \frac{1}{\sqrt{12}}(1-1+0)=0$、$\bm{p}_2, \bm{p}_3$ とも $(1,1,1)$ に直交することも確認できる）。

したがって直交行列
\[
P = \begin{pmatrix}
\dfrac{1}{\sqrt3} & \dfrac{1}{\sqrt2} & \dfrac{1}{\sqrt6} \\[6pt]
\dfrac{1}{\sqrt3} & -\dfrac{1}{\sqrt2} & \dfrac{1}{\sqrt6} \\[6pt]
\dfrac{1}{\sqrt3} & 0 & -\dfrac{2}{\sqrt6}
\end{pmatrix}
\]
により
\[
\boxed{\; P^{\mathrm{T}} A P = \begin{pmatrix} 4 & 0 & 0 \\ 0 & 1 & 0 \\ 0 & 0 & 1 \end{pmatrix} \;}
\]
（検算: $A\bm{p}_1 = (I+J)\bm{p}_1 = \bm{p}_1 + \sqrt3(1,1,1)^{\mathrm{T}} = \bm{p}_1 + 3\bm{p}_1 = 4\bm{p}_1$。また $J\bm{p}_2 = \bm{0}$（各成分の和が $0$ のため）より $A\bm{p}_2 = \bm{p}_2 + \bm{0} = 1\cdot\bm{p}_2$、同様に $A\bm{p}_3 = \bm{p}_3$。いずれも一致する。）

\clearpage
\begin{mondaiboxS}{問3}
$A = \begin{pmatrix} 1 & 2 \\ 2 & 1 \end{pmatrix}$ を対角化し、それを用いて $A^{n}$（$n$ は自然数）を求めよ。
\end{mondaiboxS}
\kaitou
固有方程式は
\[
(1-\lambda)^{2} - 4 = \lambda^{2} - 2\lambda - 3 = 0 \quad\Longrightarrow\quad \lambda = 3, -1
\]

$\lambda = 3$: $A - 3I = \begin{pmatrix} -2 & 2 \\ 2 & -2 \end{pmatrix}$ より固有ベクトル $\begin{pmatrix}1\\1\end{pmatrix}$。

$\lambda = -1$: $A + I = \begin{pmatrix} 2 & 2 \\ 2 & 2 \end{pmatrix}$ より固有ベクトル $\begin{pmatrix}1\\-1\end{pmatrix}$。

よって
\[
P = \begin{pmatrix} 1 & 1 \\ 1 & -1 \end{pmatrix}, \qquad
P^{-1} = \frac{1}{2}\begin{pmatrix} 1 & 1 \\ 1 & -1 \end{pmatrix}, \qquad
D = \begin{pmatrix} 3 & 0 \\ 0 & -1 \end{pmatrix}
\]
（$P^{-1}$ は $\det P = -2$ と余因子行列から $\dfrac{1}{-2}\begin{pmatrix}-1 & -1\\-1 & 1\end{pmatrix} = \dfrac12\begin{pmatrix}1&1\\1&-1\end{pmatrix}$ として得られる）。$A^{n} = PD^{n}P^{-1}$ を計算すると
\[
A^{n} = \begin{pmatrix} 1 & 1 \\ 1 & -1 \end{pmatrix}
\begin{pmatrix} 3^{n} & 0 \\ 0 & (-1)^{n} \end{pmatrix}
\frac{1}{2}\begin{pmatrix} 1 & 1 \\ 1 & -1 \end{pmatrix}
= \frac{1}{2}\begin{pmatrix} 3^{n} & (-1)^{n} \\ 3^{n} & -(-1)^{n} \end{pmatrix}
\begin{pmatrix} 1 & 1 \\ 1 & -1 \end{pmatrix}
\]
\[
\boxed{\; A^{n} = \frac{1}{2}\begin{pmatrix} 3^{n} + (-1)^{n} & 3^{n} - (-1)^{n} \\ 3^{n} - (-1)^{n} & 3^{n} + (-1)^{n} \end{pmatrix} \;}
\]
（検算: $n=1$ とすると $\dfrac12\begin{pmatrix}3-1 & 3+1\\3+1&3-1\end{pmatrix} = \dfrac12\begin{pmatrix}2&4\\4&2\end{pmatrix} = \begin{pmatrix}1&2\\2&1\end{pmatrix} = A$ と一致する。また $n=2$ とすると $\dfrac12\begin{pmatrix}9+1&9-1\\9-1&9+1\end{pmatrix}=\begin{pmatrix}5&4\\4&5\end{pmatrix}$ であり、直接 $A^2=\begin{pmatrix}1&2\\2&1\end{pmatrix}\begin{pmatrix}1&2\\2&1\end{pmatrix}=\begin{pmatrix}5&4\\4&5\end{pmatrix}$ と計算しても一致する。）

\end{document}
